\section{Introduction}

% Figure 1 teaser -- rendered PNG (always displays; regenerate via
% figures/make_teaser.py). Placed at the top of the Introduction (page 1).
\begin{figure*}[t]
\centering
\includegraphics[width=0.97\textwidth]{teaser.png}
\caption{\textbf{OMU-MAE in one picture.} (a)~A frozen DINOv2 ViT-B/14 turns the
camera image into a $16{\times}16$ grid of $768$-d patch features. These are
projected through the camera--LiDAR calibration $(\mathbf{P}_2,\mathbf{Tr})$ to
(b)~populate a per-voxel feature volume aligned with the LiDAR occupancy grid
($0.4$\,m, $128{\times}128{\times}32$); $85\%$ of voxels are masked (grey),
$15\%$ stay visible (coloured). (c)~A 3D CNN reconstructs occupancy and the
DINOv2 features at the masked voxels. The target encoder is frozen, so this
$2\mathrm{D}\!\rightarrow\!3\mathrm{D}$ supervision cannot collapse to a trivial
solution.}
\label{fig:teaser}
\end{figure*}


Perception models for autonomous driving are still trained
predominantly with full supervision, and the supervision they need is
uncommonly expensive: annotating LiDAR for semantic segmentation means
assigning a class to every point of every scan, a slow, specialist task
that scales poorly with fleet data volumes~\cite{semantickitti}.
Self-supervised learning (SSL) offers the standard way out. In 2D
vision, three families of pretext tasks---contrastive instance
discrimination, self-distillation such as DINO and
DINOv2~\cite{dino,dinov2}, and masked autoencoding (MAE)~\cite{mae}---now
produce representations that match or exceed supervised pretraining,
transfer to downstream tasks from a handful of labels, and in the case
of DINOv2 yield dense per-patch features with strong emergent semantics
usable by a linear classifier alone. Bringing this label efficiency to
3D driving perception is the motivation for this work.

The transfer is not immediate, because driving scenes are inherently
multi-modal: they combine camera images dense in semantics with LiDAR
point clouds dense in geometry. A useful pretrained representation must
capture \emph{both}. This raises the problem we address: \emph{how to
pretrain a 3D encoder, without any human labels and on modest compute,
so that it absorbs the semantic richness of a 2D vision foundation
model (VFM) while retaining the spatial reasoning that 3D geometry
demands}. Two failure modes make this non-trivial. Methods that learn
both modality encoders jointly must fight representational collapse and
pay a heavy compute bill; methods that copy 2D features into 3D without
a spatial pretext task never force the encoder to reason about scene
structure at all.

Existing work sits on either side of this divide, and each side
inherits one of the two limitations. Multi-modal masked autoencoders
such as UniM\textsuperscript{2}AE~\cite{unim2ae} and
NS-MAE~\cite{nsmae} bring the MAE inductive bias to driving but
reconstruct \emph{raw} signals---pixels, points, or renderings---so the
model spends capacity on photometric and measurement nuisance rather
than semantics, and their jointly trained two-branch architectures are
compute-heavy. Distillation methods such as SLidR~\cite{slidr},
ScaLR~\cite{scalr} and CleverDistiller~\cite{cleverdistiller} instead
transfer \emph{frozen VFM features} to a LiDAR encoder---a
semantically strong, collapse-free target---but they distill at every
visible position, so nothing in the objective requires the encoder to
infer missing context from surrounding structure; the contrastive
variants further operate at sparse superpoint resolution, coarser than
the dense predictions downstream tasks need. What is missing is a
method that combines the strengths of the two families: the
\emph{semantic abstraction of a frozen VFM target} with the
\emph{context-inpainting inductive bias of masked reconstruction}.

OMU-MAE (\emph{Occupancy, Multi-modal, Unified} Masked Autoencoder)
occupies exactly this position. Each scene is voxelized on a dense
grid; frozen DINOv2 patch features are projected onto the voxels
through the camera--LiDAR calibration; 85\% of voxels are masked with
Occupancy-MAE's range-aware schedule~\cite{occmae}; and a 3D CNN
encoder--decoder reconstructs both binary occupancy and the DINOv2
features at masked positions. The two design choices are each doing
identifiable work. The \emph{frozen} target cannot co-adapt to the
predictor, which rules out the constant-output collapse mode of
joint-target prediction~\cite{ijepa,lejepa} by construction---no
regularization machinery needed---while still supplying dense semantic
supervision. \emph{Masked reconstruction} forces the encoder to inpaint
missing voxel context from surrounding visible voxels, an inductive
bias that distillation at every position does not impose. Whether that
bias actually improves downstream transfer is an empirical question,
and we answer it directly: we re-implement a contrastive (SLidR-style)
and a non-masked distillation (CleverDistiller-equivalent) baseline on
identical data, architecture and evaluation protocol, so the comparison
isolates the objective itself.

\paragraph{Contributions.}
\begin{enumerate}
\item OMU-MAE, a voxel-level masked autoencoder whose reconstruction
target combines binary occupancy with frozen DINOv2~\cite{dinov2}
features projected from the paired camera image; a dual loss, proposed
here, recovers both signals at masked positions.
\item A \emph{five-way} controlled study on KITTI / SemanticKITTI
(random init / Occupancy-MAE re-impl.\ / SLidR re-impl.\ / no-mask
CleverDistiller-equivalent / OMU-MAE) under one linear-probe protocol,
isolating each design choice: the cross-modal target is worth
$+1.77$\,pp over LiDAR-only masking, the masking bias $+3.27$\,pp over
non-masked distillation, and masked reconstruction $+4.22$\,pp over
contrastive distillation of the \emph{same} frozen-VFM target.
\item A deliberately lightweight recipe---dense voxel grid, single
3D CNN branch, no interaction module---that pretrains in a few hours on
one consumer GPU, with code and pretrained weights released.
\end{enumerate}
